\section{Lifting through every boundary neighborhood}\label{sec:lifting}

We retain the data \(L,s,G,q,r,a\), the local root charts, and the maps
\(g:E\to U\), \(h:H\to U\) of Section~\ref{sec:covers}.  In particular
\(E\) is reduced Cartier, \(I=\OO_Z(-E)\) is invertible, and
\(\mathcal A_j=I^j/I^{j+1}\) has the Laurent graded frame \(u^j\) for
every integer \(j\).  The frame uses the prescribed boundary root
comparison; it is a frame on the associated graded, not an extension
of that boundary frame to \(Z\).

The finite quotients \(\mathcal T_{j,k}=I^j/I^{j+k}\) have the lifting
target in Proposition~\ref{prop:all-orders}.  The intervening
constructions supplied the split residue insertion and the projective
symbol-kernel vanishing.  We now use them to kill the connecting maps of
\eqref{eq:current-truncation-sequence} at every finite order.  The
quotients remain sheaves on the underlying space of \(E\); no map from
an infinitesimal thickening to the parameter space is required.

\subsection{The obstruction is a graded derivation}

We prove the proposition by induction on \(k\), simultaneously in all
integer degrees and all root charts.  The following description uses
only the orders strictly smaller than the current one.

\begin{lemma}\label{lem:obstruction-derivation}
Fix \(k\geq1\), and assume \eqref{eq:all-order-surjection} for every
smaller order, every integer degree, and every root chart.  The connecting
maps for the current order factor uniquely as maps of complex vector
space sheaves
\begin{equation}\label{eq:obstruction-map}
 \delta_k:g_*\mathcal A_j\longrightarrow R^1g_*\mathcal A_{j+k}
 \quad(j\in\Z).
\end{equation}
They form a degree-\(k\) derivation from the Laurent graded algebra
\(\bigoplus_jg_*\mathcal A_j\) to its graded cohomology module
\(\bigoplus_jR^1g_*\mathcal A_j\).  If \(t_1,\ldots,t_b\) are local
coordinates on \(U\) and \(F\) is holomorphic in those coordinates, then
\begin{equation}\label{eq:obstruction-chain-rule}
 \delta_k(F(t))=\sum_{i=1}^b F_{t_i}(t)\,\delta_k(t_i).
\end{equation}
These maps and formulas commute with restriction and with the ambient
power-compatible root comparisons.  Vanishing of \(\delta_k\) in all
degrees proves the current order of \eqref{eq:all-order-surjection}.
\end{lemma}
\begin{proof}
Use \eqref{eq:current-truncation-sequence} together with
\begin{equation}
 0\longrightarrow\mathcal T_{j+1,k-1}\longrightarrow\mathcal T_{j,k}
 \longrightarrow\mathcal A_j\longrightarrow0\quad(k>1).
 \label{eq:leading-truncation-sequence}
\end{equation}
All are valid for negative \(j\) because \(I\) is invertible.  Write
\(\partial_{j,k}:g_*\mathcal T_{j,k}\to R^1g_*\mathcal A_{j+k}\) for
the connecting map of the first sequence.  The shorter orders make
\(g_*\mathcal T_{j,k}\to g_*\mathcal A_j\) surjective.  For \(k>1\),
its kernel is the image of \(g_*\mathcal T_{j+1,k-1}\), by left exactness
in the second sequence.  The shorter order \(k-1\) in degree \(j+1\)
lifts that image through \(g_*\mathcal T_{j+1,k}\), which maps into
\(g_*\mathcal T_{j,k+1}\).  The connecting map vanishes on the kernel.
For \(k=1\), the leading map is the identity and its kernel is zero.
This proves the unique factorization \eqref{eq:obstruction-map}.
If it is zero, exactness of the first sequence gives the desired
epimorphism.

Here is the derivation calculation on germs.  A leading section in
degree \(j\) lifts to a section of \(\mathcal T_{j,k}\) after a base
shrink.  Choose local representatives \(x_\alpha\in I^j\) on an ambient
open cover near \(E\).  Their differences belong to \(I^{j+k}\), and
their classes modulo \(I^{j+k+1}\) represent \(\delta_k(x)\).  For a
second section \(z\) of degree \(l\), take simultaneous representatives.
On an overlap,
\[
 x_\beta z_\beta-x_\alpha z_\alpha
 =x_\alpha(z_\beta-z_\alpha)+z_\alpha(x_\beta-x_\alpha)
 +(x_\beta-x_\alpha)(z_\beta-z_\alpha).
\]
The last term lies in \(I^{j+l+2k}\subset I^{j+l+k+1}\).  The first two
terms reduce to multiplication by the leading coefficients.  Their
\v Cech classes therefore give
\[
 \delta_k(xz)=x\delta_k(z)+z\delta_k(x).
\]
This uses the ordinary multiplication action of layer sections on their
first cohomology, and gives the asserted graded derivation.

For the chain rule, choose simultaneous shorter representatives
\(G_{i,\alpha}\in\OO_Z\) of the degree-zero sections \(t_i\).  Their
differences lie in \(I^k\).  Shrink the local opens so that
\(F(G_{1,\alpha},\ldots,G_{b,\alpha})\) is defined.  Its holomorphic
Taylor difference, modulo the square of the differences, is
\(\sum_iF_{t_i}(G_\alpha)(G_{i,\beta}-G_{i,\alpha})\).
The square lies in \(I^{2k}\subset I^{k+1}\), and the coefficient
restricts to \(F_{t_i}(t)\) on \(E\).  This gives
\eqref{eq:obstruction-chain-rule} before any restriction to a generic
parameter.  It is valid when the target contains torsion.  The connecting
map and all representative calculations are natural for restrictions
and the prescribed ambient root isomorphisms.
\end{proof}

\subsection{The local differential-operator identity}

Fix the current order \(k\), and put \(m=a+k>0\).  For a local leading
section \(x\in g_*\mathcal A_{-m}\), define, in coordinates
\(t=(t_1,\ldots,t_b)\) on \(U\),
\begin{equation}\label{eq:residue-obstructions}
 c(x)=\tau_*\delta_k(x),\qquad
 e_i(x)=\tau_*\bigl(x\delta_k(t_i)\bigr)
 \quad\text{in }R^1h_*\omega_H.
\end{equation}
Both arguments of \(\tau_*\) have layer degree \(-a\), since
\(-m+k=-a\).  The map is the residue insertion of
Proposition~\ref{prop:residue-insertion}; it is injective on every
parameter germ.

Consider a holomorphic \(\psi:U'\to U\) between manifolds of the same
dimension \(b\), in coordinates \(w=(w_1,\ldots,w_b)\), satisfying
the transversality \eqref{eq:graph-transversality} for every smooth
component intersection of \(H\).  The identity map is one such map.
As in Proposition~\ref{prop:global-graph}, after a shrink about a compact
fiber the support
\[
 H^\flat=H\times_U U'\subset\widehat Z\times U'
\]
is reduced SNC, locally a divisor in a smooth complete intersection,
of codimension \(b+1\) in the product.  Its ambient projection
\(p^\flat\) is a submersion and its closed strata are proper with global
relative K\"ahler forms.  After the restrictions needed for the root
comparisons and the K\"ahler neighborhood, apply
Lemma~\ref{lem:finite-components} anew to the smooth closed strata of
the graph cut over \(U'\), and use its final open neighborhood.  Properness
survives this base change, and the lemma gives finitely many globally
smooth irreducible component labels and finitely many smooth proper
connected pieces of each closed intersection.  The existing coordinates
and relative forms restrict to that open.  Set
\[
 K^\flat=\mathcal H^{b+1}_{[H^\flat]}(\omega_{\widehat Z\times U'}),
 \qquad M^\flat=\mathcal H^1p^\flat_+K^\flat.
\]
Proposition~\ref{prop:filtered-snc} applies with no additional
codimension shift and gives
\(F_0M^\flat=R^1h^\flat_*\omega_{H^\flat}\hookrightarrow M^\flat\).

The adjunction identity \eqref{eq:graph-adjunction} supplies the factor
\(\omega_{U'}\otimes\psi^*\omega_U^{-1}\).  We denote its wedge frame
by \(dw/dt\); this notation is a ratio of canonical frames, not an
inverse Jacobian.  Pull the \v Cech classes representing \(c(x),e_i(x)\)
to \(H^\flat\) and tensor with this frame.  Write their images in
\(F_0M^\flat\) as \(\widetilde c,\widetilde e_i\).  These are the natural
pullbacks of these classes; no base-change isomorphism is asserted for
a ramified \(\psi\).

\begin{lemma}[The adjugate identity]\label{lem:local-symbol}
Let \(J=(\partial\psi_i/\partial w_j)\) and
\(J^\#=\operatorname{adj}(J)\).  Under the preceding assumptions, the
following identity holds in the right module \(M^\flat\):
\begin{equation}\label{eq:adjugate-operator}
 \widetilde c\,\det J+
 \sum_{i,j}(\widetilde e_i\,\partial_{w_j})J^\#_{ji}=0.
\end{equation}
In particular,
\begin{equation}\label{eq:local-symbol-zero}
 \sigma\left(\sum_{j,i}J^\#_{ji}\widetilde e_i\otimes\partial_{w_j}\right)=0
 \quad\text{in }\gr^F_1M^\flat.
\end{equation}
For \(\psi=\id\), after the natural graph identification, the exact
identity is
\begin{equation}\label{eq:identity-operator}
 c(x)+\sum_i e_i(x)\partial_{t_i}=0\quad\text{in }M^\flat.
\end{equation}
When \(b=0\), the same construction without graph equations gives
\(c(x)=0\) in its degree-one direct-image module.
\end{lemma}
\begin{proof}
All statements are on germs, so choose a common parameter shrink and
simultaneous shorter lifts of the finitely many sections \(x,t_i\).
Choose local ambient representatives near \(E\), indexed by \(\alpha\),
with
\begin{equation}\label{eq:shorter-representatives}
 x_\alpha\in I^{-a-k},\quad x_\beta-x_\alpha\in I^{-a},\qquad
 G_{i,\alpha}\in\OO_Z,\quad G_{i,\beta}-G_{i,\alpha}\in I^k.
\end{equation}
They exist by the smaller orders in Lemma~\ref{lem:obstruction-derivation};
for \(k=1\) the shorter truncation is already the layer.  Pull them to
\(\widehat Z\), suppressing pullback symbols, and put
\[
 \eta_\alpha=x_\alpha\tau,\qquad
 Q_{i,\alpha}=\psi_i(w)-G_{i,\alpha},\qquad
 P_\alpha=\prod_i Q_{i,\alpha}.
\]
On \(H\) the \(G_{i,\alpha}\) restrict to the coordinate functions of
\(h\).  The \(Q_i\)'s and a reduced equation of \(H\) are therefore
the regular graph-support equations.  The logarithmic insertion and
\(D_E=\rho^*E\geq H\) give the precise pole bounds
\begin{align}
 \eta_\alpha&\in\omega_{\widehat Z}(H+kD_E),\nonumber\\
 \eta_\beta-\eta_\alpha,\quad
 \eta_\alpha(G_{i,\beta}-G_{i,\alpha})
 &\in\omega_{\widehat Z}(H),\label{eq:fraction-pole-bounds}\\
 \eta_\alpha(G_{i,\beta}-G_{i,\alpha})(G_{l,\beta}-G_{l,\alpha}),
 \quad(\eta_\beta-\eta_\alpha)(G_{i,\beta}-G_{i,\alpha})
 &\in\omega_{\widehat Z}(H-kD_E)\subset\omega_{\widehat Z}.
 \nonumber
\end{align}
The last inclusion uses \(k\geq1\).  Thus the quadratic and cross
difference terms have no pole on the first-factor support \(H\).

In algebraic local cohomology consider the local sections
\[
 S_\alpha=\left[\frac{\eta_\alpha\wedge dw}{P_\alpha}\right]
 \quad\text{of }K^\flat.
\]
The \(H\)-pole is already in \(\eta_\alpha\).  These are finite-pole
generalized fractions with a fixed ordering of the graph equations.
To expand their differences, first quotient the localization along
\(H\) by forms regular there, and then localize for the \(Q_i\)'s.
Each polar class is killed by a power of a reduced equation of \(H\),
and every \(G_{i,\beta}-G_{i,\alpha}\) is a multiple of that equation.
The change from \(Q_{i,\alpha}\) to
\(Q_{i,\beta}=Q_{i,\alpha}-(G_{i,\beta}-G_{i,\alpha})\) therefore has
a finite geometric expansion on each polar class.  By
\eqref{eq:fraction-pole-bounds} only its linear terms survive.  This
proves the exact generalized-fraction equality on an overlap
\begin{equation}\label{eq:fraction-difference}
 S_\beta-S_\alpha=
 \left[\frac{(\eta_\beta-\eta_\alpha)\wedge dw}{P_\alpha}\right]
 +\sum_i\left[
 \frac{\eta_\alpha(G_{i,\beta}-G_{i,\alpha})\wedge dw}
 {Q_{i,\alpha}P_\alpha}\right].
\end{equation}
There is no infinite series or division by a Jacobian in this identity.

Let \(C\) be the first \v Cech cochain in
\eqref{eq:fraction-difference}.  Let \(E_i\) have the numerator of its
\(i\)th summand but only the denominator \(P_\alpha\), and let
\(E_i^{(l)}\) denote that cochain with the additional denominator
\(Q_{l,\alpha}\).  The simple fractions \(C,E_i\) are cocycles: their
reductions are respectively the cocycles for \(\delta_k(x)\) and
\(x\delta_k(t_i)\); on a triple overlap the possible change of the
representative of \(x\) is a cross term with no \(H\)-pole by
\eqref{eq:fraction-pole-bounds}.  The description of the residue insertion on representatives and
complete-intersection adjunction identify their classes with
\(\widetilde c,\widetilde e_i\) in \(F_0M^\flat\).  More explicitly,
the residue of a log form \(\eta\) on \(H\), pulled with \(dw/dt\), is
represented by \([\eta\wedge dw/P_\alpha]\); the determinant of the
graph conormal gives exactly that frame.  Thus this identification holds
also at ramification.  The cochains \(E_i^{(l)}\) need not individually
be cocycles.

The numerator of \(E_i\) is independent of \(w\).  The right canonical
action is minus differentiation of the coefficient of the volume form,
so at the cochain level
\begin{equation}\label{eq:right-fraction-action}
 E_i\partial_{w_j}=\sum_l J_{lj}E_i^{(l)}.
\end{equation}
Write \(\check d\) for the \v Cech differential.  Equation
\eqref{eq:fraction-difference} is
\(\check dS=C+\sum_iE_i^{(i)}\).  Multiply it by \(\det J\) and use
\(JJ^\#=(\det J)\id\) together with
\eqref{eq:right-fraction-action}.  The result, still at the cochain
level, is
\begin{equation}\label{eq:cochain-adjugate}
 \check d(S\det J)=C\det J+
 \sum_{i,j}(E_i\partial_{w_j})J^\#_{ji}.
\end{equation}
The placement of \(J^\#_{ji}\) after the right derivative is part of
this exact formula.

These cochains lie in the end term of the relative right de Rham complex
for the product projection, with first-factor coordinates held fixed.
That term has no outgoing relative differential, so a \v Cech coboundary
there is a total coboundary in direct image.  The action in
\eqref{eq:right-fraction-action} induces the base right action there.
Passing \eqref{eq:cochain-adjugate} to \(M^\flat\) gives
\eqref{eq:adjugate-operator}.  This uses the natural map from these
\v Cech classes to direct image, and does not require that all direct-image
classes be computed by this cover.

The term \(\widetilde c\det J\) is in \(F_0\).  Moving a holomorphic
function past a right vector field changes an operator by order zero.
Taking \(\gr^F_1\) of \eqref{eq:adjugate-operator} gives
\eqref{eq:local-symbol-zero}.  With unchanged coordinates \(J=\id\),
the exact formula itself gives \eqref{eq:identity-operator}.  If \(b=0\),
there are no \(Q_i\); the undivided difference of the \(\eta_\alpha\)
in the support module of \(H\) is the cochain for \(c(x)\) and is a total
coboundary.  This proves the last assertion as well.
\end{proof}

\subsection{One global symbol kills the current obstruction}

\begin{proof}[Proof of Proposition~\ref{prop:all-orders}]
Assume the shorter orders and use Lemma~\ref{lem:obstruction-derivation}.
First suppose \(b>0\), and fix an arbitrary point \(t_*\in T=\PP^b\).
Proposition~\ref{prop:global-graph} gives the coordinate-power map
\(\psi:T'=\PP^b\to T\), the line \(R=\OO_{T'}(1)\) with
\(R^r\simeq\psi^*\OO_T(1)\), and a single compact SNC graph support
\(H'\subset\mathscr W\) with submersive projection \(p':\mathscr W\to T'\).
It is unbranched over \(t_*\).  Put
\[
 K'=\mathcal H^{b+1}_{[H']}(\omega_{\mathscr W}),\qquad
 M'=\mathcal H^1p'_+K'.
\]
The global graph proposition verifies every hypothesis of
Proposition~\ref{prop:filtered-snc}: the support has pure dimension
\(\dim V-1\), globally smooth finitely many components, and proper
smooth closed intersections with one ambient K\"ahler form.  Thus
\(F_0M'=R^1h'_*\omega_{H'}\hookrightarrow M'\), in the absolute
canonical convention, and Corollary~\ref{cor:symbol-kernel} applies.

Take the local leading section \(x=u^{-m}\).  On a local graph chart
choose a frame \(\bar p\) of \(R\) with \(\bar p^r=\psi^*\ell\) for a
downstairs frame \(\ell\).  The local formula
\begin{equation}\label{eq:global-symbol-local}
 \bar p^m\longmapsto
 \sum_{j,i}J^\#_{ji}\widetilde e_i(u^{-m})\otimes\partial_{w_j}
\end{equation}
defines a global morphism
\begin{equation}\label{eq:global-symbol-map}
 R^m\longrightarrow F_0M'\otimes T_{T'}.
\end{equation}
We verify the transition, including on the branch locus.

Let \(t'=t'(t)\) and \(w'=w'(w)\) be changes of coordinates, and put
\(A_t=\partial t'/\partial t\), \(B_w=\partial w'/\partial w\).  These
matrices are invertible.  If \(\bar p'=\lambda\bar p\), then
\(\lambda^r=\psi^*\mu\) for the downstairs unit \(\mu\) changing the
frame of \(\OO_T(1)\).  Locally even at a branch point, \(\lambda\)
descends as a holomorphic unit: take an \(r\)th root of \(\mu\) on a
small downstairs neighborhood, and observe that the remaining ratio has
\(r\)th power one and hence is locally constant.  Denote the resulting
downstairs unit also by \(\lambda\).

Compare the local root pairs by sending their new boundary frame to
\(\lambda p_1\).  Lemma~\ref{lem:root-neighborhood} extends this
power-compatible comparison to the ambient germs.  Its pullback agrees
with the transition of the global root near the graph, by the uniqueness
and properness argument in Proposition~\ref{prop:global-graph}.  It
therefore identifies both full covers, their tautological forms, and
the fixed functorial principalizations.  Under this identification
\(u'^{-m}=\lambda^m u^{-m}\).  In the definition of \(e_i(x)\), the
factor \(x\) multiplies outside the derivation.  Hence no derivative of
\(\lambda\) occurs.  Naturality and the chain rule
\eqref{eq:obstruction-chain-rule} transform the downstairs column of
\(e_i\)'s by \(\lambda^m A_t\).

The absolute graph factor changes by
\(dw'/dt'=(\det B_w/\det A_t)\,dw/dt\).  Consequently, for the pulled
columns and Jacobians,
\begin{align}
 \widetilde e'&=\frac{\det B_w}{\det A_t}\lambda^m A_t\widetilde e,
 &J'&=A_tJB_w^{-1},\label{eq:symbol-transition-one}\\
 (J')^\#&=\frac{\det A_t}{\det B_w}\,B_wJ^\#A_t^{-1},
 &(J')^\#\widetilde e'&=\lambda^m B_wJ^\#\widetilde e.
 \label{eq:symbol-transition-two}
\end{align}
The adjugate identity holds for every \(J\): it follows for invertible
\(J\) from the inverse formula and then for all \(J\) as a polynomial
identity.  Only the coordinate changes \(A_t,B_w\) are inverted.
Since the vector basis changes by \(B_w^{-1}\), the last equality is
exactly the transition for the image of the frame \(\bar p^m\) in
\eqref{eq:global-symbol-local}.  This proves \eqref{eq:global-symbol-map}
on all of \(T'\).  At ramification the classes are the natural \v Cech
pullbacks under the resolved ambient germ comparisons; a flat base-change
isomorphism for \(R^1h_*\omega_H\) is not used.

By Lemma~\ref{lem:local-symbol}, the image of
\eqref{eq:global-symbol-map} lies in \(\ker\sigma\).  The line
\(R^m=\OO_{\PP^b}(a+k)\) is ample because \(a+k>0\).  Corollary
\ref{cor:symbol-kernel} makes the global map zero.  At a point above
\(t_*\), the map \(\psi\) is locally biholomorphic and \(J^\#\) is
invertible.  The ambient product comparison then identifies the graph
family and its \v Cech classes with the downstairs family.  Thus the
zero map implies \(e_i(u^{-m})=0\) as a downstairs germ near \(t_*\),
including a class supported at that parameter.  The point \(t_*\) was
arbitrary, and the root comparisons allow any chart at it.  It follows
that these classes vanish locally everywhere.

The injection \(\tau_*\) and the invertible layer frame \(u^{-m}\)
now imply \(\delta_k(t_i)=0\) on every chart: multiplication by that
frame is an isomorphism from the layer \(\mathcal A_k\) to
\(\mathcal A_{-a}\), including on first direct images.  For any local
\(x\) of degree \(-m\), all \(e_i(x)\) vanish.  The exact unchanged
coordinate identity \eqref{eq:identity-operator} gives \(c(x)=0\) in
\(M^\flat\).  The injection \eqref{eq:lowest-direct-image} brings this
back to \(R^1h_*\omega_H\), and the residue injection gives
\(\delta_k(x)=0\).  When \(b=0\), Lemma~\ref{lem:local-symbol} gives
\(c(x)=0\) directly, and the same two injections give this conclusion;
there are no parameter coordinates to treat.

In particular \(\delta_k(u^{-m})=0\).  The derivation rule for the
invertible Laurent frame gives
\[
 0=\delta_k(u^{-m})=-m u^{-m-1}\delta_k(u),
\]
so \(\delta_k(u)=0\).  For any local \(F\in g_*\OO_E\), the section
\(Fu^{-m}\) also has degree \(-m\), and
\[
 0=\delta_k(Fu^{-m})=u^{-m}\delta_k(F).
\]
Thus \(\delta_k(F)=0\).  Every local graded section is \(Fu^j\), so
\(\delta_k\) is zero in every integer degree.  Lemma
\ref{lem:obstruction-derivation} closes the current order.  Induction
from \(k=1\) proves Proposition~\ref{prop:all-orders}.
\end{proof}

\subsection{Invariant layers and the growth of sections}

Finite lifting is now available in every degree.  To obtain ambient
sections, we pass to cyclic invariants and count the resulting finite
quotients downstairs.  Translation by the Cartier divisor \(G\) will
make the same finitely many coherent layers recur with increasing
positive twists on \(T\).  Those twists supply the section growth and
bound the higher-cohomology loss.

\begin{proof}[Completion of Theorem~\ref{thm:supported-lifting}]
Define global divisorial subsheaves of meromorphic functions on the
normal space \(V\) by
\begin{equation}\label{eq:downstairs-ideals}
 J_j=\OO_V\left(-\ceil{jG/r}\right)\quad(j\in\Z).
\end{equation}
They agree in each root chart with \((\pi_*I^j)^{\mu_r}\).  Indeed, at
a prime over \(S_i\), an invariant meromorphic function of downstairs
order \(v_i\) has order \((r/d_i)v_i\).  Membership in \(I^j\) is
equivalent to
\[
 v_i\geq\ceil{jd_i/r}.
\]
At all other primes it is regular.  The full invariant meromorphic
algebra is the downstairs algebra, and normality extends the
codimension-one test.  Finite coherent pushforward and averaging by
\(\mu_r\) are exact.  This also proves coherence of the divisorial
sheaves in \eqref{eq:downstairs-ideals} and of their quotients locally,
and their valuation descriptions glue globally.

Since \(0<d_i\leq r\), the sheaf \(J_1=\OO_V(-S)\) is the ideal of the
reduced divisor \(S\).  The same coefficient inequality gives
\(J_1J_j\subset J_{j+1}\) for every \(j\).  Hence
\[
 \mathcal B_j=J_j/J_{j+1}
\]
is a coherent \(\OO_S\)-module, and \(\mathcal B_0=\OO_S\).
Taking invariants of Proposition~\ref{prop:all-orders} preserves its
epimorphisms: lift an invariant local section upstairs and average the
lift.  For \(l<n\), define the sheaf of complex vector spaces
\[
 \mathscr F_{l,n}=f_*(J_l/J_n)\quad\text{on }T,
 \qquad F_j=f_*\mathcal B_j.
\]
Set also \(\mathscr F_{n,n}=0\).
The quotients \(J_l/J_n\) are considered on the underlying space of
\(S\), just as before.  The invariant lifting surjections give exact
sequences of complex vector space sheaves
\begin{equation}\label{eq:downstairs-filtration}
 0\longrightarrow\mathscr F_{j+1,n}\longrightarrow\mathscr F_{j,n}
 \longrightarrow F_j\longrightarrow0\quad(j<n).
\end{equation}
In particular \(\mathscr F_{l,n}\) has a finite filtration with the
full quotients \(F_j\), \(l\leq j<n\).  Each \(F_j\) is a coherent
\(\OO_T\)-module by properness of \(f\); coherence as an \(\OO_T\)-module
is not asserted for \(\mathscr F_{l,n}\).

Cartier translation by \(G\) in \eqref{eq:downstairs-ideals}, together
with the actual line identity \(\OO_V(G)=L\) defined by \(s\), gives
\[
 J_{j-r}=J_j\otimes\OO_V(G),\qquad
 \mathcal B_{j-r}=\mathcal B_j\otimes L|_S.
\]
Projection formula and \eqref{eq:boundary-map} consequently give
\begin{equation}\label{eq:periodic-layers}
 F_{j-r}=F_j\otimes\OO_T(1),\qquad F_0=f_*\OO_S\ne0.
\end{equation}
This is an identity of the actual holomorphic lines and sheaves.

For \(N>0\), each \(-Nr\leq j<0\) is uniquely \(j=l-tr\) with
\(0\leq l<r\) and \(1\leq t\leq N\).  By
\eqref{eq:periodic-layers}, \(F_j=F_l(t)\).  Analytic GAGA and Serre's
coherent finiteness and vanishing theorems
\cite[Section~3, no.~12, Theorems~1 and~3]{Serre1956GAGA},
\cite[Section~66, Theorems~1 and~2]{Serre1955FAC} imply, for every
\(i>0\),
\begin{equation}\label{eq:higher-layer-bound}
 \sum_{j=-Nr}^{-1}h^i(T,F_j)
 \leq C_i:=\sum_{l=0}^{r-1}\sum_{t\geq1}h^i(T,F_l(t))<\infty.
\end{equation}
These constants are independent of \(N\), and they are zero for
\(i>\dim T\).  The long exact sequences of the finite filtration
\eqref{eq:downstairs-filtration} give the same bounds for
\(h^i(T,\mathscr F_{-Nr,0})\), as well as finiteness and vanishing above
that fixed dimension.  Serre vanishing is being applied only to the
coherent quotients \(F_l(t)\).

There is at least one section of \(F_0(t)\) for every \(t\geq1\).
Choose a point of the nonempty \(S\) and a section of \(\OO_T(t)\)
nonzero at its image; its pullback is nonzero there.  Thus
\(\sum_{j=-Nr}^{-1}h^0(T,F_j)\geq N\).  Additivity of the finite Euler
characteristics in \eqref{eq:downstairs-filtration}, and the uniform
higher bounds for both the quotients and \(\mathscr F\), give a constant
\(C\) independent of \(N\) with
\begin{equation}\label{eq:thickening-growth}
 h^0(T,\mathscr F_{-Nr,0})\geq N-C.
\end{equation}
For example one can take \(C=2\sum_{i=1}^{\dim T}C_i\).

Finally, \(J_{-Nr}=\OO_V(NG)\) and \(J_0=\OO_V\).  Degree-zero direct
image on the underlying support gives
\[
 H^0(T,\mathscr F_{-Nr,0})=
 H^0(V,\OO_V(NG)/\OO_V).
\]
The exact sequence on \(V\) loses at most the fixed finite dimension
\(h^1(V,\OO_V)\) when lifting these quotient sections to
\(H^0(V,\OO_V(NG))\); finiteness follows from proper coherent direct
image to a point.  Equation \eqref{eq:thickening-growth} therefore makes
these section spaces unbounded.  For some \(N\) the line \(\OO_V(NG)\)
has two independent sections.  Their quotient is a nonconstant meromorphic
function on the irreducible normal \(V\), so its Iitaka dimension is at
least one.  Since \(\OO_V(G)=L=\OO_V(qA)\) is the actual adjoint
identity, this proves \(\kappa(V,A)\geq1\), as claimed.
\end{proof}
